\section{AI Agent：规划、推理与工具使用}

\subsection{什么是 AI Agent}

讲师指出，``AI Agent'' 这个词在不同人口中含义差异极大，从简单的 LLM API 调用到复杂的多 LLM 协作系统都被冠以``Agent''之名。在讲师看来，Agent 最核心的两个能力维度是：

\begin{enumerate}
    \item \textbf{规划与推理}（Planning \& Reasoning）：能够分解复杂任务、制定计划、反思中间结果
    \item \textbf{工具使用}（Tool Use）：能够调用外部 API（搜索引擎、计算器、数据库等）获取信息或执行操作
\end{enumerate}

\subsection{ReAct：推理与行动的融合}

讲师介绍了 ReAct（Reasoning + Acting）论文，这是 Agent 领域最具影响力的奠基性工作之一。

在 ReAct 之前，研究者要么只让模型做推理（reasoning traces），要么只让模型执行动作（action traces）。ReAct 将两者结合为一个混合框架。

\textbf{具体对比}：

\textit{任务}：判断 ``Reign Over Me is an American film made in 2010'' 是否为真。

\begin{table}[H]
\centering
\begin{tabular}{p{6cm}p{6cm}}
\toprule
\textbf{Chain-of-Thought（仅推理）} & \textbf{ReAct（推理 + 行动）} \\
\midrule
Thought: It is an American film... & Thought: I need to search for Reign Over Me \\
Thought: It was made in 2010... & Action: Search{[}Reign Over Me{]} \\
Answer: \textbf{True} {\color{red}（错误）} & Observation: American film, released \textbf{2007} \\
 & Thought: It was made in 2007, not 2010 \\
 & Answer: \textbf{False} {\color{green!60!black}（正确）} \\
\bottomrule
\end{tabular}
\caption{Chain-of-Thought vs ReAct 的对比}
\end{table}

纯 CoT 方法因为只靠模型``内部知识''推理，错误地认为该电影是 2010 年拍摄的（实际是 2007 年）。ReAct 通过搜索外部信息源获得了正确的发行年份，从而给出正确答案。

\begin{importantbox}{ReAct 框架的循环结构}
ReAct 的核心是 Thought $\rightarrow$ Action $\rightarrow$ Observation 的循环：
\begin{enumerate}
    \item \textbf{Thought}：模型推理当前状态，决定下一步该做什么
    \item \textbf{Action}：模型执行一个动作（如搜索、调用 API）
    \item \textbf{Observation}：环境返回动作的结果
    \item 重复上述循环，直到模型认为已有足够信息给出最终答案
\end{enumerate}
这个框架已成为现代 Agent 系统（如 LangChain、AutoGPT 等）的基础架构。
\end{importantbox}

讲师还展示了另一个 ReAct 的例子：一个文本冒险游戏，任务是``将胡椒瓶放到抽屉上''。

\begin{itemize}
    \item \textbf{仅行动}模式：Agent 走到水槽 1，试图拿不存在的胡椒瓶，陷入循环
    \item \textbf{ReAct 模式}：Agent 推理``水槽 1 没有胡椒瓶，需要去别处找''，成功导航并完成任务
\end{itemize}

\subsection{Toolformer：教 LLM 使用工具}

讲师介绍的第二篇奠基性论文是 Meta 的 Toolformer，它教会 LLM 在生成文本时自主调用外部 API。

Toolformer 支持的工具类型包括：
\begin{itemize}
    \item \textbf{QA 工具}：回答事实性问题
    \item \textbf{Calculator}：执行精确计算
    \item \textbf{Machine Translation}：翻译
    \item \textbf{Wikipedia Search}：查找百科知识
\end{itemize}

\begin{knowledgebox}{Toolformer 的训练方法——三步流水线}
\begin{enumerate}
    \item \textbf{标注生成}：给模型少量示例（few-shot），让模型在文本中标注``在这里应该调用某个 API''
    \item \textbf{API 执行}：实际调用这些 API，获取返回结果
    \item \textbf{有效性过滤}（这是最巧妙的一步）：比较``包含 API 结果''和``不包含 API 结果''两种情况下的训练 loss。只保留那些\textbf{显著降低了 loss} 的 API 调用作为训练数据
\end{enumerate}
例如，``Pittsburgh was known as the {[}QA: Steel City{]}'' 这个 API 调用有用（模型原本不太确定 Pittsburgh 的昵称），而 ``the country that Pittsburgh is in is {[}QA: United States{]}'' 这个调用没用（模型本来就知道），因此后者被过滤掉。
\end{knowledgebox}

\subsection{本章小结}

AI Agent 的两大核心能力——规划推理和工具使用——由 ReAct 和 Toolformer 两篇奠基性论文分别开拓。ReAct 将推理（Thought）和行动（Action）融合为一个循环框架，使 Agent 能够动态调整策略；Toolformer 则教会模型在何时以及如何调用外部工具，并通过巧妙的过滤机制确保只学习有价值的工具调用模式。

